\documentclass[../../main.tex]{subfiles}% 注意这里的文件路径不能用 ./main.tex ,否则用latexmk编译子文件会报错
\graphicspath{{\subfix{./image/}}} % 指定图片目录,后续可以直接使用图片文件名
% 注意这里的文件路径不能用 ../../image/ ,否则用latexmk编译子文件会报错

% 例如:
% \begin{figure}[H]
% \centering
% \includegraphics[scale=0.3]{图.png}
% \caption{}
% \label{figure:图}
% \end{figure}
% 注意:上述\label{}一定要放在\caption{}之后,否则引用图片序号会只会显示??.

\begin{document}

\section{保长对应和保角对应}

\begin{definition}
设正则参数曲面$S_1:\boldsymbol{r}=\boldsymbol{r}_1(u_1,v_1),\,(u_1,v_1)\in D_1$，$S_2:\boldsymbol{r}=\boldsymbol{r}_2(u_2,v_2),\,(u_2,v_2)\in D_2$，其中$D_1,D_2\subset\mathbb{E}^2$。若映射$\sigma:D_1\to D_2$表示为
\begin{align}
u_2=f(u_1,v_1),\quad v_2=g(u_1,v_1),\label{eq:surfmap}
\end{align}
则$\sigma$诱导曲面间的映射（仍记为$\sigma$）：$S_1\to S_2$，满足$\sigma(\boldsymbol{r}_1(u_1,v_1))=\boldsymbol{r}_2(f(u_1,v_1),g(u_1,v_1))$。
若$f,g$为$3$次以上连续可微函数，则称$\sigma:S_1\to S_2$为\textbf{3阶光滑映射}。

对任意$p\in S_1$，设$C_1:u_1=u_1(t),\,v_1=v_1(t)$为$S_1$上过$p$的曲线，其在$\sigma$下的像为$C_2:u_2=f(u_1(t),v_1(t)),\,v_2=g(u_1(t),v_1(t))$。定义\textbf{切映射}$\sigma_{*p}:T_pS_1\to T_{\sigma(p)}S_2$：
\begin{align}
\sigma_{*p}\left(\left.\frac{\mathrm{d}\boldsymbol{r}_1}{\mathrm{d}t}\right|_{t=t_0}\right)=\left.\frac{\mathrm{d}\boldsymbol{r}_2}{\mathrm{d}t}\right|_{t=t_0}.\label{eq:tangentmap}
\end{align}

设$\varphi=A(u_2,v_2)(\mathrm{d}u_2)^2+2B(u_2,v_2)\mathrm{d}u_2\mathrm{d}v_2+C(u_2,v_2)(\mathrm{d}v_2)^2$为$S_2$上的二次微分式，通过映射\eqref{eq:surfmap}代入得到$S_1$上的二次微分式$\sigma^*\varphi$，称$\sigma^*\varphi$为$\varphi$经$\sigma$的\textbf{拉回微分式}，满足
\begin{align*}
\sigma^*\varphi=(\mathrm{d}u_1,\mathrm{d}v_1)J\begin{pmatrix}A&B\\B&C\end{pmatrix}J^\mathrm{T}\begin{pmatrix}\mathrm{d}u_1\\\mathrm{d}v_1\end{pmatrix},
\end{align*}
其中$J=\begin{pmatrix}\frac{\partial u_2}{\partial u_1}&\frac{\partial v_2}{\partial u_1}\\[4pt]\frac{\partial u_2}{\partial v_1}&\frac{\partial v_2}{\partial v_1}\end{pmatrix}$为映射的Jacobi矩阵。
\end{definition}

\begin{proposition}
切映射$\sigma_{*p}:T_pS_1\to T_{\sigma(p)}S_2$为线性映射；$\sigma_{*p}$为线性同构当且仅当Jacobi行列式$\left.\frac{\partial(u_2,v_2)}{\partial(u_1,v_1)}\right|_p\neq0$。
\end{proposition}
\begin{proof}

对任意切向量$\boldsymbol{X}=a\frac{\partial\boldsymbol{r}_1}{\partial u_1}+b\frac{\partial\boldsymbol{r}_1}{\partial v_1}\in T_pS_1$，由链式法则
\begin{align*}
\sigma_{*p}(\boldsymbol{X})&=a\left(\frac{\partial\boldsymbol{r}_2}{\partial u_2}\frac{\partial f}{\partial u_1}+\frac{\partial\boldsymbol{r}_2}{\partial v_2}\frac{\partial g}{\partial u_1}\right)+b\left(\frac{\partial\boldsymbol{r}_2}{\partial u_2}\frac{\partial f}{\partial v_1}+\frac{\partial\boldsymbol{r}_2}{\partial v_2}\frac{\partial g}{\partial v_1}\right)\\
&=\left(a\frac{\partial u_2}{\partial u_1}+b\frac{\partial u_2}{\partial v_1}\right)\frac{\partial\boldsymbol{r}_2}{\partial u_2}+\left(a\frac{\partial v_2}{\partial u_1}+b\frac{\partial v_2}{\partial v_1}\right)\frac{\partial\boldsymbol{r}_2}{\partial v_2},
\end{align*}
故$\sigma_{*p}$保持线性运算，为线性映射。
线性同构等价于$\frac{\partial\boldsymbol{r}_2}{\partial u_2},\frac{\partial\boldsymbol{r}_2}{\partial v_2}$的系数矩阵可逆，即Jacobi行列式非零。

\end{proof}

\begin{theorem}
设$\sigma:S_1\to S_2$为3阶光滑映射，若$\sigma_{*p}:T_pS_1\to T_{\sigma(p)}S_2$为同构，则存在$p$的邻域$U_1\subset S_1$、$\sigma(p)$的邻域$U_2\subset S_2$，及参数系$(u_1,v_1),(u_2,v_2)$，使得$\sigma|_{U_1}$满足$u_2=u_1,\,v_2=v_1$。
\end{theorem}
\begin{proof}

由Jacobi行列式$\left.\frac{\partial(u_2,v_2)}{\partial(u_1,v_1)}\right|_p\neq0$，根据隐函数定理，$u_2=f(u_1,v_1),v_2=g(u_1,v_1)$可作为$S_1$在$p$邻域内的容许参数变换，此时$\sigma$在参数区域上为恒同映射。

\end{proof}

\begin{definition}
设$\sigma:S_1\to S_2$为3阶光滑映射，若对任意$p\in S_1$及任意$\boldsymbol{X}\in T_pS_1$，满足$|\sigma_{*p}(\boldsymbol{X})|=|\boldsymbol{X}|$，则称$\sigma$为\textbf{保长对应}。
\end{definition}

\begin{proposition}
保长对应保持切向量的内积不变，即对任意$\boldsymbol{X},\boldsymbol{Y}\in T_pS_1$，有$\sigma_{*p}(\boldsymbol{X})\cdot\sigma_{*p}(\boldsymbol{Y})=\boldsymbol{X}\cdot\boldsymbol{Y}$。
\end{proposition}
\begin{proof}

由向量内积恒等式$\boldsymbol{a}\cdot\boldsymbol{b}=\frac{1}{2}\big(|\boldsymbol{a}+\boldsymbol{b}|^2-|\boldsymbol{a}|^2-|\boldsymbol{b}|^2\big)$，结合保长对应保持长度不变，直接可得内积不变。

\end{proof}

\begin{theorem}
设$S_1,S_2$的第一基本形式分别为$\mathrm{I}_1,\mathrm{I}_2$，则$\sigma:S_1\to S_2$为保长对应的充要条件是$\sigma^*\mathrm{I}_2=\mathrm{I}_1$，即
\begin{align}
\begin{pmatrix}E_1&F_1\\F_1&G_1\end{pmatrix}=J\begin{pmatrix}E_2&F_2\\F_2&G_2\end{pmatrix}J^\mathrm{T},\label{eq:isomcond}
\end{align}
其中$E_i,F_i,G_i$为$S_i$的第一类基本量，$J$为映射的Jacobi矩阵。
\end{theorem}
\begin{proof}

$\sigma$为保长对应$\iff|\mathrm{d}\boldsymbol{r}_1|^2=|\sigma_{*p}(\mathrm{d}\boldsymbol{r}_1)|^2$，
而$|\sigma_{*p}(\mathrm{d}\boldsymbol{r}_1)|^2=\sigma^*(|\mathrm{d}\boldsymbol{r}_2|^2)=\sigma^*\mathrm{I}_2$，
故等价于$\mathrm{I}_1=\sigma^*\mathrm{I}_2$，即矩阵形式\eqref{eq:isomcond}。

\end{proof}

\begin{theorem}
存在$S_1$到$S_2$的保长对应的充要条件是：可选取公共参数系$(u,v)$，使得$E_1=E_2,\,F_1=F_2,\,G_1=G_2$。
\end{theorem}
\begin{proof}

必要性：若$\sigma$为保长对应，则$\sigma_{*p}$为同构，由前述定理可选取适用参数系，此时$\sigma$为恒同参数对应，代入保长条件得第一类基本量相等。
充分性：若公共参数系下第一类基本量相等，则$\sigma$取参数恒等映射，满足$\sigma^*\mathrm{I}_2=\mathrm{I}_1$，故为保长对应。

\end{proof}

\begin{example}
证明螺旋面$\boldsymbol{r}=(u\cos v,u\sin v,u+v)$与旋转双曲面$\boldsymbol{r}=(\rho\cos\theta,\rho\sin\theta,\sqrt{\rho^2-1})\,(\rho\geqslant1,0\leqslant\theta\leqslant2\pi)$之间可建立保长对应。
\begin{solution}

计算螺旋面的第一基本形式：
\begin{align*}
\mathrm{I}=2(\mathrm{d}u)^2+2\mathrm{d}u\mathrm{d}v+(u^2+1)(\mathrm{d}v)^2.
\end{align*}
旋转双曲面的第一基本形式：
\begin{align*}
\tilde{\mathrm{I}}=\frac{2\rho^2-1}{\rho^2-1}(\mathrm{d}\rho)^2+\rho^2(\mathrm{d}\theta)^2.
\end{align*}
对螺旋面作参数变换$\tilde{u}=u,\,\tilde{v}=\arctan u+v$，代入化简得
\begin{align*}
\mathrm{I}=\frac{2\tilde{u}^2+1}{\tilde{u}^2+1}(\mathrm{d}\tilde{u})^2+(\tilde{u}^2+1)(\mathrm{d}\tilde{v})^2.
\end{align*}
令$\rho=\sqrt{\tilde{u}^2+1},\,\theta=\tilde{v}$，则$\mathrm{I}=\tilde{\mathrm{I}}$，故映射$\rho=\sqrt{u^2+1},\,\theta=\arctan u+v$为保长对应。

\end{solution}
\end{example}


\begin{definition}
设$\sigma:S_1\to S_2$为微分同胚（$\sigma$与逆映射均为3阶光滑），若对任意$p\in S_1$及任意$\boldsymbol{X},\boldsymbol{Y}\in T_pS_1$，满足$\angle(\sigma_{*p}(\boldsymbol{X}),\sigma_{*p}(\boldsymbol{Y}))=\angle(\boldsymbol{X},\boldsymbol{Y})$，则称$\sigma$为\textbf{保角对应}。
\end{definition}

\begin{theorem}
$\sigma:S_1\to S_2$为保角对应的充要条件是：存在$S_1$上正的连续函数$\lambda$，使得$\sigma^*\mathrm{I}_2=\lambda^2\mathrm{I}_1$。
此时可选取公共参数系$(u,v)$，使得
\begin{align}
E_2=\lambda^2E_1,\quad F_2=\lambda^2F_1,\quad G_2=\lambda^2G_1.\label{eq:confcond}
\end{align}
\end{theorem}
\begin{proof}

保角对应即切向量夹角不变，等价于内积成比例、长度成比例，即$\sigma^*\mathrm{I}_2=\lambda^2\mathrm{I}_1$。
选取适用参数系后，直接得到基本量成比例的条件\eqref{eq:confcond}。

\end{proof}

\begin{theorem}
设正则参数曲面$S_1,S_2$的第一基本形式分别为$\mathrm{I}_1,\mathrm{I}_2$，则映射$\sigma:S_1\to S_2$为\textbf{保角对应}的充分必要条件是：存在$S_1$上的正连续函数$\lambda$，使得
\begin{align}
\sigma^*\mathrm{I}_2 = \lambda^2\mathrm{I}_1.
\label{eq:conformalcond}
\end{align}
特别地，若$\sigma$为保角对应，则可在$S_1,S_2$上取相同参数记号$(u,v)$，使得两曲面的第一类基本量成比例，即
\begin{align}
\begin{cases}
E_2(u,v)=\lambda^2(u,v)E_1(u,v),\\
F_2(u,v)=\lambda^2(u,v)F_1(u,v),\\
G_2(u,v)=\lambda^2(u,v)G_1(u,v).
\end{cases}
\label{eq:basicalpropconformal}
\end{align}
\end{theorem}

\begin{proof}

{\heiti 充分性}:
若\eqref{eq:conformalcond}成立，任取$p\in S_1$，$\bm X,\bm Y\in T_pS_1$，由切向量夹角公式
\begin{align}
\cos\angle(\bm X,\bm Y)=\frac{\mathrm{I}_1(\bm X,\bm Y)}{\sqrt{\mathrm{I}_1(\bm X,\bm X)}\sqrt{\mathrm{I}_1(\bm Y,\bm Y)}},\label{eq:anglecos1}\\
\cos\angle(\sigma_{*p}(\bm X),\sigma_{*p}(\bm Y))=\frac{\sigma^*\mathrm{I}_2(\bm X,\bm Y)}{\sqrt{\sigma^*\mathrm{I}_2(\bm X,\bm X)}\sqrt{\sigma^*\mathrm{I}_2(\bm Y,\bm Y)}}.\label{eq:anglecos2}
\end{align}
结合$\sigma^*\mathrm{I}_2(\cdot,\cdot)=\lambda^2\mathrm{I}_1(\cdot,\cdot)$，代入可得
$$\cos\angle(\sigma_{*p}(\bm X),\sigma_{*p}(\bm Y))=\cos\angle(\bm X,\bm Y),$$
故$\sigma$保角。

{\heiti 必要性}:
若$\sigma$为保角对应，则$\sigma$为一一对应，可取相同参数系$(u,v)$使得$\sigma(\bm r_1(u,v))=\bm r_2(u,v)$，此时
$$\sigma_*\left(\frac{\partial\bm r_1}{\partial u}\right)=\frac{\partial\bm r_2}{\partial u},\quad \sigma_*\left(\frac{\partial\bm r_1}{\partial v}\right)=\frac{\partial\bm r_2}{\partial v}.$$
由保角性，参数曲线切向量夹角、切向量和的夹角均保持不变，代入夹角公式得
\begin{align*}
\frac{F_1}{\sqrt{E_1G_1}}=\frac{F_2}{\sqrt{E_2G_2}},\quad
\frac{E_2+F_2}{\sqrt{E_2+2F_2+G_2}\sqrt{E_2}}=\frac{E_1+F_1}{\sqrt{E_1+2F_1+G_1}\sqrt{E_1}},\quad
\frac{F_2+G_2}{\sqrt{E_2+2F_2+G_2}\sqrt{G_2}}=\frac{F_1+G_1}{\sqrt{E_1+2F_1+G_1}\sqrt{G_1}}.
\end{align*}
联立化简可得$E_1G_2=E_2G_1$，进一步推得
$$\frac{F_2}{F_1}=\frac{E_2}{E_1}=\frac{G_2}{G_1}=\lambda^2,$$
即基本量成比例，式\eqref{eq:conformalcond}成立.

\end{proof}

\begin{theorem}
对任意正则参数曲面$S$上一点$p$，存在$p$的邻域$U\subset S$，使得$U$可与平面开区域建立保角对应；进一步，任意两个正则参数曲面在局部均可建立保角对应。
\end{theorem}

\begin{remark}
该定理为曲面论深刻结论：解析曲面情形可利用复解析函数与一次微分式积分因子证明；光滑曲面情形证明较复杂；$C^k(k\geqslant2)$类曲面仍成立。
\end{remark}

\begin{proof}

设$S:\bm r=\bm r(u,v)$为解析曲面，由正交参数系存在性，可取正交参数系使得$F=0$，此时第一基本形式为
\begin{align*}
\mathrm{I}=E(\mathrm{d}u)^2+G(\mathrm{d}v)^2=\big(\sqrt{E}\mathrm{d}u+\sqrt{-1}\sqrt{G}\mathrm{d}v\big)\big(\sqrt{E}\mathrm{d}u-\sqrt{-1}\sqrt{G}\mathrm{d}v\big).
\end{align*}
令复一次微分式
\begin{align}
\omega=\sqrt{E}\mathrm{d}u+\sqrt{-1}\sqrt{G}\mathrm{d}v.
\label{eq:complexdiff}
\end{align}
存在非零复解析函数$\lambda(u,v)$，使得$\lambda\omega=\mathrm{d}z(u,v)$，其中$z= x+\sqrt{-1}y$为复解析函数。代入得
\begin{align*}
\mathrm{I}=\omega\overline{\omega}=\frac{1}{|\lambda|^2}\big(\mathrm{d}x^2+\mathrm{d}y^2\big),
\end{align*}
故映射$(u,v)\mapsto(x(u,v),y(u,v))$为$U$到平面的保角对应.

\end{proof}

\begin{definition}
曲面$S$上使得第一基本形式满足
$$\mathrm{I}=\frac{1}{|\lambda|^2}\big(\mathrm{d}x^2+\mathrm{d}y^2\big)$$
的参数系$(x,y)$，称为$S$的\textbf{等温参数}。
\end{definition}

\begin{example}
建立半径为$a$的球面与同半径圆柱面之间的保角对应（Mercator投影）。
\end{example}

\begin{solution}

球面$S$与圆柱面$\tilde S$的参数方程分别为
\begin{align*}
\bm r&=\big(a\cos v\cos u,\,a\cos v\sin u,\,a\sin v\big),\\
\tilde{\bm r}&=\big(a\cos\tilde u,\,a\sin\tilde u,\,a\tilde v\big).
\end{align*}
计算第一基本形式：
\begin{align*}
\mathrm{I}&=a^2\cos^2v(\mathrm{d}u)^2+a^2(\mathrm{d}v)^2=a^2\cos^2v\left((\mathrm{d}u)^2+\frac{1}{\cos^2v}(\mathrm{d}v)^2\right),\\
\tilde{\mathrm{I}}&=a^2(\mathrm{d}\tilde u)^2+a^2(\mathrm{d}\tilde v)^2.
\end{align*}
取参数变换
\begin{align*}
\tilde u=u,\quad \tilde v=\int_{0}^{v}\frac{\mathrm{d}t}{\cos t}=\ln\left|\tan\left(\frac{v}{2}+\frac{\pi}{4}\right)\right|,
\end{align*}
则$\mathrm{I}=a^2\cos^2v\tilde{\mathrm{I}}$，满足保角对应条件，该映射即为\textbf{Mercator投影}。

圆柱面与平面保长，故球面可通过该投影与平面局部保角对应，常用于地图绘制.

\end{solution}

\begin{example}
\begin{enumerate}[(1)]
\item 证明: 在悬链面
\begin{align*}
\boldsymbol{r}=(a\cosh t\cos\theta,a\cosh t\sin\theta,at),\quad -\infty<t<\infty,\ 0\leqslant\theta\leqslant2\pi
\end{align*}
和正螺面
\begin{align*}
\boldsymbol{r}=(v\cos u,v\sin u,au),\quad 0\leqslant u\leqslant2\pi,\ -\infty<v<\infty
\end{align*}
之间存在保长对应.

\item 证明: 曲面
\begin{align*}
\boldsymbol{r}=(a(\cos u+\cos v),a(\sin u+\sin v),b(u+v))
\end{align*}
能够和一个旋转面建立保长对应.

\item 证明: 平面到它自身的任意一个保长对应必定是平面上的一个刚体运动,或刚体运动和关于一条直线的反射的合成.

\item 试建立旋转面
\begin{align*}
\boldsymbol{r}=(f(u)\cos v,f(u)\sin v,g(u))
\end{align*}
和平面的保角对应.

\item 试建立第(2)题中的曲面和平面的保角对应.
\end{enumerate}
\end{example}



\end{document}